\section{Rounded edges and the cost of a flat disk reference}
\label{sec:ed03-rounded-edges}

\begin{theorem}[Birth bodies and edge incidence]
\label{thm:ed03-birth-edge-incidence}
\uses{thm:ed03-native-foot-charts}
\lean{SmoothPoincare.IsolatedZeroHandle.birthBodyChart_component, SmoothPoincare.IsolatedZeroHandle.exists_transportedBirthBody, SmoothPoincare.IsolatedOneHandle.traceComponentMap_fibers, SmoothPoincare.IsolatedOneHandle.levelFeet_disjoint_of_bodyFoot_ne, SmoothPoincare.IsolatedOneHandle.not_simplyConnected_of_same_bodyFoot, SmoothPoincare.IsolatedMiddleAttachment.traceComponentMap_bijective}
\leanok
At a selected birth at $p$, the radius-one closed four-ball in its rescaled
Morse chart is the whole sublevel component of $p$ at $f(p)+\varepsilon$.
The chart extends across its boundary and the ball and its marked sphere
transport together across the following regular band.
For a selected one-handle, the map on body components induced by the lower
insertion and the trace diffeomorphism is onto. Its only identifications are
the component labels of the two native feet. Distinct feet in the same body
preclude a simply connected outgoing body. The corresponding component map
for a selected two-handle is bijective.
\end{theorem}
\begin{proof}
\leanok
The birth chart has height $f(p)+\varepsilon |x|^2$. Its closed unit ball
is connected, compact and open relative to the selected sublevel. To compute
the component map, collapse the component labels of the attaching region
in the adjunction. The handle is connected. For a loop edge the joined-feet
obstruction applies to the actual trace. For an index-two attachment the circle
of normal-disk centers is connected, so it carries a single incoming label.
\end{proof}

\begin{theorem}[Rounded comparison at both native rims]
\label{thm:ed03-rounded-native-rims}
\uses{thm:ed03-birth-edge-incidence,thm:ed03-marked-edge-feet,thm:ed03-native-height-collar}
\lean{SmoothPoincare.IsolatedOneHandle.edgeBlockChart_nativeFoot, SmoothPoincare.IsolatedOneHandle.edgeBlockChart_afterFootNormalization, SmoothPoincare.IsolatedOneHandle.edgeBlockChart_signedOutgoing, SmoothPoincare.IsolatedOneHandle.roundedRimToOutgoing_signed_inverse, SmoothPoincare.IsolatedOneHandle.exists_roundedOutgoingCollar_nearNativeRims, SmoothPoincare.IsolatedOneHandle.exists_outgoingSignedBandTransport, SmoothPoincare.IsolatedOneHandle.exists_roundedRimTransport_toNextCut}
\leanok
The four-dimensional block chart restricts to both entire native attaching
disks with their parameters. Supplied foot normalizations retain their
individual parities. For a rounding on the Morse core with radius exceeding
$\sqrt{2\varepsilon+2r^2}$, its radial comparison $T$ has a smooth inverse
near both full rims. If $\Phi$ is the actual outgoing product chart obtained
in the construction, there are an open set $U$ containing both rim images and
$0<w\le\varepsilon/2$ on which the rounded collar satisfies
\[
 \Psi=T^{-1}\circ\Phi,\qquad
 \Psi^{-1}=\Phi^{-1}\circ T,\qquad
 f_{\rm rounded}(\Psi(z,t))=f(p)+t
 \quad(z\in U,\ |t|<w).
\]
Let $B_1$ be the collar at the next regular cut, $d$ the regular level
comparison, and $H$ the band comparison based on this same $\Phi$. On
a common positive substrip the maps satisfy
\[
 H(T(\Psi(z,t)))=B_1(dz,t),\qquad
 T^{-1}(H^{-1}(B_1(dz,t)))=\Psi(z,t).
\]
These equations retain every tube parameter supplied in the outgoing level,
and on the rim strip they retain the original label and both normal directions.
\end{theorem}
\begin{proof}
\leanok
At fixed axial square $s$, add $2\varepsilon+s-\operatorname{smoothCap}
(\varepsilon,r,\delta,s)$ to the squared transverse radius. Subtraction
recovers the starting point on the inverse domain. Both native rims lie
in the common open core. Compactness of $S^0\times S^2$ gives a uniform
two-sided interval for the pullback of the outgoing chart. At the next cut
conjugate the level transport times the identity by the two signed collars.
\end{proof}
A global edge contraction still needs the two incoming body markings, the
remaining tubes and compatible maps on their open overlaps with the rounded
block. A product extension chosen from two collars may shear tangentially
relative to another extension that agrees on the zero section.

\begin{theorem}[Realized reference on a selected width]
\label{thm:ed03-realized-disk-reference}
\uses{prop:ed03-physical-marked-jets,prop:ed03-selected-quadratic-jets}
\lean{SmoothPoincare.exists_selectedEuclideanReferenceSplice, SmoothPoincare.exists_selectedEuclideanSecondPartialError_on_physicalTrim, SmoothPoincare.selectedEuclideanSplice_flat_near_innerDisk, SmoothPoincare.eventually_balancedFlatReference, SmoothPoincare.selectedReference_endpoint_saving}
\leanok
Let $Q\subset\operatorname{int} B$ be the nested marked disks on one physical
source, and $W_n$ its selected trim. For every positive tolerance all sufficiently
large $n$ give a single width for which the four original trim masses have
sum below tolerance and $Q\cap W_n$ is empty. On that width one can splice
a smooth positive definite Euclidean reference in the original coordinates.
The time-one metric is flat on $B\setminus W_n$, including $Q$. All exterior
metric germs on $(B\cup W_n)^c$ agree throughout the path.
\end{theorem}
\begin{proof}
\leanok
Extend the reference from the buffered chart by a smooth cutoff. Convexity of
positive definite forms gives the path. The cutoff is constant on the inner
and outer open regions. The physical jet estimate uses the original measure
and the same width. The unit seam can cross the transition. A smaller closed disk
has an open flat neighborhood at positive $n$.
\end{proof}

\begin{theorem}[Signed price of a positive flat replacement]
\label{thm:ed03-mandatory-flat-cost}
\uses{thm:ed03-realized-disk-reference}
\lean{SmoothPoincare.balancedFlatReference_cost, SmoothPoincare.balancedFlatReference_trim_cost, SmoothPoincare.eventually_balancedFlatReference_cost, SmoothPoincare.eventually_balancedFlatReference_trim_charge_pos, SmoothPoincare.selectedReference_charge_eq_of_germ}
\leanok
Put $q=2^{-k}$, $I=\inf E_+$ over positive Yamabe classes on the fixed
oriented smooth type, $e=E_+(g)$ and $R=(B\cup W)^c$.
If the marked exterior division at depth $k$ has error $\epsilon$ and $h$
is an attained Yamabe-positive metric flat on $B\setminus W$, then for
each $s\in\{+,-\}$,
\[
 \lambda_h^s(W)+\lambda_h^s(R)-\lambda_g^s(R)
 > I-q\bigl(e-\lambda_g^s(Q)\bigr)-\epsilon .
\]
Matching exterior germs and frames put the full lower bound on $W$.
When $e_i\to I$ and $\epsilon_i\to0$, for every $\eta>0$ the left side
exceeds $(1-q)I-\eta$ for both signs at all sufficiently late indices,
for every such positive endpoint. If $I>0$ and $k>0$, an endpoint preserving
the full exterior charge must eventually have strictly positive trim charge
in each sign.
\end{theorem}
\begin{proof}
\leanok
The flat portion contributes zero, so the two new charges sum to
$E_s(h)\ge I$. Signature zero gives this inequality in both chiralities.
The full incoming exterior controls $\lambda_g^s(R)$ by the marked division.
Subtract and pass to the stated eventual inequality, keeping the finite core
charge nonnegative. This argument requires no minimizing metric at $I$.
\end{proof}

\begin{proposition}[Exact saving on a matched flat trim]
\label{prop:ed03-native-matched-flat}
\uses{thm:ed03-realized-disk-reference,thm:ed03-mandatory-flat-cost}
\lean{SmoothPoincare.exists_nativeConstantMetric_at, SmoothPoincare.exists_selectedMatchedFlatDisk, SmoothPoincare.selectedMatchedFlatDisk_trim_charge, SmoothPoincare.selectedMatchedFlatDisk_exact_saving, SmoothPoincare.selectedSplice_exterior_chiralMeasure_restrict_eq}
\leanok
Suppose the incoming metric in the given marking agrees with a constant
positive Gram matrix $L^*L$ on an open neighborhood of the complete closed
trim. The path to that flat reference has zero trim curvature at every time
and matches the original unit seam. Its endpoint is flat on the whole disk.
For an attained positive endpoint $h$ with the stated exterior frame agreement,
\[
 E_s(h)=\lambda_g^s((\operatorname{int}B)^c),\qquad
 E_+(g)-E_+(h)=\lambda_g^s(B)\le e-I\qquad(s=+,-).
\]
Equality of restricted measures for independent exterior frames is also
proved when their orthogonal action on two-forms is supplied with Hodge
projection and curvature-conjugacy identities.
\end{proposition}
\begin{proof}
\leanok
Pull back the Euclidean metric by a constant linear isomorphism to define the
reference tensor in the native parameters. The interpolation is constant on the
transition. Flatness of the endpoint and equality on the complete exterior
give the two charge identities. Membership of the positive endpoint in the
fixed-type energy set gives the upper bound on its saving.
\end{proof}

\begin{proposition}[Euler flux on the supplying flow]
\label{prop:ed03-spatial-euler-flux}
\uses{thm:ed03-modified-euler-transport}
\lean{SmoothPoincare.metricFrameEulerPfaffian_eq_chargeDensity, SmoothPoincare.metricFrameEulerCovariantDerivative_bianchi, SmoothPoincare.frameEulerPfaffian_hasDerivAt_of_bianchi, SmoothPoincare.ricciFlow_frameEulerTransgression_norm, SmoothPoincare.eulerCurvatureCoefficient_hasDerivAt_skewJet, SmoothPoincare.ricciFlowEulerPfaffian_fixedBase_hasDerivAt, SmoothPoincare.eulerSpatialFlux_integral_bound, SmoothPoincare.modifiedEuler_transport_of_spatialFlux}
\leanok
In an oriented orthonormal frame the normalized scalar Euler density is
twice the Pfaffian coefficient in the frame convention used here. The
transgression coefficients obey $|T|\le144|\mathrm{Rm}|\,|\nabla\mathrm{Ric}|$
in the proved frame norm. On an actual smooth Ricci solution the coefficient
variation is computed with the four base vectors fixed and endomorphism-frame
velocities $\mathrm{Ric}^{\sharp}e_a$ at the tested time. Suppose a smooth slice current $X$ realizes
$m'(s)=\int\chi\operatorname{div}_{g(s)}X_s\,dV_s$ in the open time interval.
With the corresponding current bound, slice curvature norm at most $A$,
integrated squared derivative norm at most $L^2$, and fixed smooth cutoff
$|d\chi|_s\le\delta(\tau/s)^b$, one has
\[
 |m(\tau)-m(0)|\le144\,\delta A L\sqrt{\tau/(1-2b)}
 \qquad (b<1/2,\ \tau>0)
\]
under the stated endpoint continuity and flux measurability.
\end{proposition}
\begin{proof}
\leanok
Apply the differential Bianchi identity to the frame product rule and the
Ricci equation to the orthogonal connection variation. Spatial integration
by parts and Cauchy--Schwarz on each supplying slice, followed by weighted
Cauchy--Schwarz in time, give the estimate. The current, the two endpoint
moments and every slice volume belong to the same evolution.
\end{proof}
